\documentclass[11pt]{article}
\newcommand{\doctitle}{pmsimstats-ng Codebase Overview}
\newcommand{\shorttitle}{Codebase Overview}
\newcommand{\docdate}{2026-04-09}
% pmsimstats-ng standard document preamble
% Usage: % pmsimstats-ng standard document preamble
% Usage: % pmsimstats-ng standard document preamble
% Usage: \input{pmsimstats-preamble} after \documentclass[11pt]{article}
%
% Documents must define before \input:
%   \newcommand{\doctitle}{...}       Full document title
%   \newcommand{\shorttitle}{...}     Short title for header
%   \newcommand{\docdate}{YYYY-MM-DD} Document date

\usepackage[margin=1in]{geometry}
\usepackage{amsmath,amssymb}
\usepackage[T1]{fontenc}
\usepackage{lmodern}
\usepackage{microtype}
\usepackage{graphicx}
\graphicspath{{figures/}}
\usepackage{booktabs}
\usepackage{longtable}
\usepackage{array}
\usepackage{hyperref}
\usepackage{xcolor}
\usepackage{fancyhdr}
\usepackage{parskip}
\usepackage{caption}
\usepackage{enumitem}
\usepackage{verbatim}
\usepackage{listings}

\lstset{
  language=R,
  basicstyle=\small\ttfamily,
  keywordstyle=\color{blue!70!black},
  commentstyle=\color{green!50!black},
  stringstyle=\color{red!60!black},
  breaklines=true,
  frame=single,
  framerule=0.5pt,
  rulecolor=\color{gray!40},
  backgroundcolor=\color{gray!5},
  xleftmargin=1em,
  framexleftmargin=1em,
  columns=flexible
}

\hypersetup{
  colorlinks=true,
  linkcolor=blue!60!black,
  citecolor=green!50!black,
  urlcolor=blue!60!black
}

\pagestyle{fancy}
\fancyhf{}
\lhead{\footnotesize pmsimstats-ng Technical Report}
\rhead{\footnotesize \shorttitle}
\rfoot{\footnotesize \thepage}
\renewcommand{\headrulewidth}{0.4pt}

\title{\doctitle}
\author{pmsimstats team}
\date{\docdate}
 after \documentclass[11pt]{article}
%
% Documents must define before \input:
%   \newcommand{\doctitle}{...}       Full document title
%   \newcommand{\shorttitle}{...}     Short title for header
%   \newcommand{\docdate}{YYYY-MM-DD} Document date

\usepackage[margin=1in]{geometry}
\usepackage{amsmath,amssymb}
\usepackage[T1]{fontenc}
\usepackage{lmodern}
\usepackage{microtype}
\usepackage{graphicx}
\graphicspath{{figures/}}
\usepackage{booktabs}
\usepackage{longtable}
\usepackage{array}
\usepackage{hyperref}
\usepackage{xcolor}
\usepackage{fancyhdr}
\usepackage{parskip}
\usepackage{caption}
\usepackage{enumitem}
\usepackage{verbatim}
\usepackage{listings}

\lstset{
  language=R,
  basicstyle=\small\ttfamily,
  keywordstyle=\color{blue!70!black},
  commentstyle=\color{green!50!black},
  stringstyle=\color{red!60!black},
  breaklines=true,
  frame=single,
  framerule=0.5pt,
  rulecolor=\color{gray!40},
  backgroundcolor=\color{gray!5},
  xleftmargin=1em,
  framexleftmargin=1em,
  columns=flexible
}

\hypersetup{
  colorlinks=true,
  linkcolor=blue!60!black,
  citecolor=green!50!black,
  urlcolor=blue!60!black
}

\pagestyle{fancy}
\fancyhf{}
\lhead{\footnotesize pmsimstats-ng Technical Report}
\rhead{\footnotesize \shorttitle}
\rfoot{\footnotesize \thepage}
\renewcommand{\headrulewidth}{0.4pt}

\title{\doctitle}
\author{pmsimstats team}
\date{\docdate}
 after \documentclass[11pt]{article}
%
% Documents must define before \input:
%   \newcommand{\doctitle}{...}       Full document title
%   \newcommand{\shorttitle}{...}     Short title for header
%   \newcommand{\docdate}{YYYY-MM-DD} Document date

\usepackage[margin=1in]{geometry}
\usepackage{amsmath,amssymb}
\usepackage[T1]{fontenc}
\usepackage{lmodern}
\usepackage{microtype}
\usepackage{graphicx}
\graphicspath{{figures/}}
\usepackage{booktabs}
\usepackage{longtable}
\usepackage{array}
\usepackage{hyperref}
\usepackage{xcolor}
\usepackage{fancyhdr}
\usepackage{parskip}
\usepackage{caption}
\usepackage{enumitem}
\usepackage{verbatim}
\usepackage{listings}

\lstset{
  language=R,
  basicstyle=\small\ttfamily,
  keywordstyle=\color{blue!70!black},
  commentstyle=\color{green!50!black},
  stringstyle=\color{red!60!black},
  breaklines=true,
  frame=single,
  framerule=0.5pt,
  rulecolor=\color{gray!40},
  backgroundcolor=\color{gray!5},
  xleftmargin=1em,
  framexleftmargin=1em,
  columns=flexible
}

\hypersetup{
  colorlinks=true,
  linkcolor=blue!60!black,
  citecolor=green!50!black,
  urlcolor=blue!60!black
}

\pagestyle{fancy}
\fancyhf{}
\lhead{\footnotesize pmsimstats-ng Technical Report}
\rhead{\footnotesize \shorttitle}
\rfoot{\footnotesize \thepage}
\renewcommand{\headrulewidth}{0.4pt}

\title{\doctitle}
\author{pmsimstats team}
\date{\docdate}


\begin{document}
\maketitle
\thispagestyle{fancy}
\tableofcontents

The \texttt{pmsimstats-ng} R package implements Monte Carlo
simulation-based power analysis for aggregated N-of-1 clinical
trial designs, comparing four trial designs on their ability to
detect a biomarker-treatment interaction. The package descends
from the publication code of Hendrickson et al.\ (2020),
\textit{Frontiers in Digital Health} 2:13. Since the merge on
2026-04-08, the package supports two data-generating-process
architectures: Architecture~B (multivariate normal differential
correlation, the original Hendrickson approach) and
Architecture~A (direct mean moderation), selectable via the
\texttt{dgp\_architecture} parameter. The dual-architecture
distinction is the subject of the white paper
\texttt{02-dgp-mean-moderation-vs-mvn.tex}.

\section{Core R Package Files (\texttt{R/})}

\subsection{generateData.R}

The data-generating process (DGP). Three exported functions:

\begin{itemize}
  \item \textbf{\texttt{generateData()}} -- Draws simulated
    participant data from a multivariate normal distribution.
    Accepts a pre-built covariance matrix
    (\texttt{cached\_sigma}) or builds one via
    \texttt{buildSigma()}. Uses a pre-computed Cholesky factor
    for efficient MVN draws (\texttt{L \%*\% Z + mu}). Computes
    outcome columns via vectorized \texttt{data.table}
    operations. The \texttt{dgp\_architecture} parameter
    selects between Architecture~B (\texttt{"mvn"}, default)
    and Architecture~A (\texttt{"mean\_moderation"}); under
    Architecture~A, the BM-BR correlation block is omitted
    from $\Sigma$ and the biomarker-treatment interaction is
    applied as an additive post-draw step
    $\text{BR}_{it} \mathrel{+}= c_{bm} \cdot B_i^{*} \cdot
    \mathbf{1}[\text{on drug}] \cdot \sigma_{\text{BR}}$,
    where $B_i^{*}$ is the standardized biomarker.

  \item \textbf{\texttt{buildSigma()}} -- Constructs the 26x26
    covariance matrix for a single trial path. Also accepts
    \texttt{dgp\_architecture}: under Architecture~A the BM-BR
    correlation block is set to zero. Implements:
    \begin{itemize}
      \item AR(1) within-factor autocorrelation: $\rho^{|t_i - t_j|}$
      \item Cross-factor correlations (same-time and different-time
        with AR(1) decay)
      \item Biomarker-BR correlation with exponential decay during
        off-drug periods: $c_{bm} \cdot \exp(-\lambda_{cor} \cdot t_{sd})$
        (Architecture~B only)
      \item Auto-derives $\lambda_{cor} = \ln(2) / \texttt{carryover\_t1half}$
      \item Positive definiteness check with verbose warning
      \item Pre-computes Cholesky factor for caching
    \end{itemize}

  \item \textbf{\texttt{validateParameterGrid()}} -- Tests all
    parameter combinations for positive definiteness before
    simulation begins, reporting failures upfront. Note: this
    function currently builds sigma matrices with Architecture~B's
    BM-BR correlation block regardless of the simulation's actual
    \texttt{dgp\_architecture} setting, so it may report spurious
    PD failures when running Architecture~A with high
    \texttt{c.bm}. The data generation itself is unaffected; the
    warning can be safely ignored under Architecture~A. See doc
    \texttt{19} for the full discussion.
\end{itemize}

\subsection{generateSimulatedResults.R}

The simulation loop. Single exported function:

\begin{itemize}
  \item \textbf{\texttt{generateSimulatedResults()}} -- Orchestrates the full
    parameter sweep:
    \begin{enumerate}
      \item PD pre-validation of all parameter combinations
      \item Sigma matrix caching (builds all unique matrices once)
      \item Parameter grid expansion across designs, response
        params, baseline params, and model params
      \item Per-parameter-set worker function
        (\texttt{run\_one\_paramset}) that generates data, fits LME,
        applies censoring, and collects results
      \item Parallel execution via \texttt{furrr::future\_map()} with
        closure-based function passing to workers
      \item Progressive save support for long runs
    \end{enumerate}

    Key parameters:
    \begin{itemize}
      \item \texttt{dgp\_architecture}: \texttt{"mvn"} (default)
        or \texttt{"mean\_moderation"}.
      \item \texttt{lambda\_cor}: correlation decay rate;
        default \texttt{NA} (auto-derived from
        \texttt{carryover\_t1half}).
      \item \texttt{n\_cores}: parallelism; default
        auto-detect.
    \end{itemize}
\end{itemize}

\subsection{lme\_analysis.R}

The analysis model. Single exported function:

\begin{itemize}
  \item \textbf{\texttt{lme\_analysis()}} -- Fits a linear mixed-effects model
    to simulated or actual trial data:
    \begin{itemize}
      \item Uses \texttt{nlme::lme} with \texttt{corCAR1(form = \~{}t|ptID)} for
        continuous-time AR(1) residual correlation
      \item Dynamic formula selection: \texttt{Sx \~{} bm + t + Dbc +
        bm:Dbc} when within-subject drug variation exists;
        \texttt{Sx \~{} bm + t + bm:t} for open-label
      \item \texttt{Dbc} is a continuous drug indicator (1 when on drug,
        exponential decay when off drug)
      \item Graceful handling of rank-deficient models (returns NA)
      \item Fallback to \texttt{lme} without \texttt{corCAR1} if optimization
        fails
      \item Extracts interaction coefficient from \texttt{nlme} tTable
    \end{itemize}
\end{itemize}

\subsection{carryover\_analysis.R}

Extended analysis functions for actual trial data:

\begin{itemize}
  \item \textbf{\texttt{characterize\_carryover()}} -- Sweeps candidate
    carryover half-lives (default 0.1 to 8 weeks), fits
    \texttt{nlme::lme} at each, compares AIC/BIC to identify the
    best-fitting half-life. Returns drug effect, interaction,
    AR(1) phi at each half-life.

  \item \textbf{\texttt{analyze\_trial\_extended()}} -- Extracts both the main
    drug effect (\texttt{Dbc}) and biomarker interaction (\texttt{bm:Dbc})
    with 95\% CIs. Computes variance components, predicted
    drug effect as a function of biomarker value, and
    clinical decision threshold.

  \item \textbf{\texttt{print\_carryover\_summary()}} / \textbf{\texttt{print\_trial\_summary()}}
    -- Clinical reporting summaries.
\end{itemize}

\subsection{buildtrialdesign.R}

Trial design specification:

\begin{itemize}
  \item \textbf{\texttt{buildtrialdesign()}} -- Constructs trial paths from
    intuitive inputs (timepoints, expectancies, on-drug
    vectors). Computes derived timing variables: \texttt{tod}
    (time on drug), \texttt{tsd} (time since discontinuation),
    \texttt{tpb} (time in positive belief).
\end{itemize}

\subsection{utilities.R}

Helper functions:

\begin{itemize}
  \item \textbf{\texttt{modgompertz()}} -- Modified Gompertz function for
    response trajectories. Passes through origin, asymptotes
    at \texttt{maxr}.
  \item \textbf{\texttt{cumulative()}} -- Converts interval durations to
    cumulative time.
  \item \textbf{\texttt{reknitsimresults()}} -- Recombines progressive save
    chunks.
\end{itemize}

\subsection{censordata.R}

Dropout simulation:

\begin{itemize}
  \item \textbf{\texttt{censordata()}} -- Two-component dropout: time-dependent
    (flat rate) and outcome-dependent (biased by symptom
    change). Implements last-observation-carried-forward
    censoring.
\end{itemize}

\subsection{plottingfunctions.R}

Visualization:

\begin{itemize}
  \item \textbf{\texttt{PlotModelingResults()}} -- Power heatmaps across
    4-dimensional parameter grids with \texttt{ggplot2 geom\_tile}
    and \texttt{facet\_grid}.
  \item \textbf{\texttt{plotfactortrajectories()}} -- Factor trajectory plots
    by trial path.
\end{itemize}

\section{Implementation Collections (\texttt{implementations/})}
\label{sec:impl}

In addition to the active package code in \texttt{R/}, the
repository contains three parallel implementations of the
simulation framework, each supporting different DGP
architectures and coding styles. The active \texttt{R/}
directory mirrors \texttt{implementations/original-extended/}.

\begin{table}[h]
\centering
\begin{tabular}{lll}
\toprule
Collection & Architectures & Style \\
\midrule
\texttt{original} & B (MVN) only & data.table \\
\texttt{original-extended} & A + B & data.table \\
\texttt{tidyverse} & A + B & tidyverse \\
\bottomrule
\end{tabular}
\end{table}

\begin{itemize}
  \item \textbf{\texttt{original}} -- Corrected Hendrickson
    publication code, restricted to Architecture~B. Useful as
    a backward-compatibility baseline and for testing the
    Hendrickson et al.\ (2020) specification in isolation. Has
    no \texttt{dgp\_architecture} parameter.

  \item \textbf{\texttt{original-extended}} -- Production
    implementation. Adds the \texttt{dgp\_architecture}
    parameter to all relevant functions, supporting both
    Architecture~A (\texttt{"mean\_moderation"}) and
    Architecture~B (\texttt{"mvn"}, default). The active
    \texttt{R/} package directory matches this collection
    file-for-file.

  \item \textbf{\texttt{tidyverse}} -- Complete tidyverse
    reimplementation. All routines live in a single
    \texttt{functions.R} file and use snake\_case names
    (e.g., \texttt{generate\_data},
    \texttt{build\_sigma\_matrix}, \texttt{lme\_analysis}).
    Supports both architectures via the same
    \texttt{dgp\_architecture} parameter.
\end{itemize}

Each collection has its own \texttt{README.md} documenting
parameter usage, validation results, and any
collection-specific behavior. All three accept the same core
parameters (\texttt{c.bm}, \texttt{carryover\_t1half},
\texttt{lambda\_cor}); only \texttt{original-extended} and
\texttt{tidyverse} accept the \texttt{dgp\_architecture}
parameter.

\section{Analysis Scripts}

\subsection{Figure 4 (\texttt{analysis/figure4/})}

\begin{table}[h]
\centering
\begin{tabular}{ll}
\toprule
Script & Purpose \\
\midrule
\texttt{01-generate-power-sweep.R} & Figure 4: current code, sequential \\
\texttt{02-generate-power-sweep-original.R} & Figure 4: original commit code \\
\texttt{03-generate-power-by-commit.R} & Figure 4: parameterized by commit \\
\texttt{04-generate-power-head.R} & Figure 4: current code, parallel, c.bm=\{0,0.25,0.45\} \\
\texttt{05-pd-diagnostics.R} & PD failure rate instrumentation \\
\texttt{06-plot-power-heatmaps.R} & Heatmap plotting with provenance \\
\bottomrule
\end{tabular}
\end{table}

\subsection{Figure 5 (\texttt{analysis/figure5/})}

\begin{table}[h]
\centering
\begin{tabular}{ll}
\toprule
Script & Purpose \\
\midrule
\texttt{01-generate-parameter-sensitivity.R} & Figure 5: revised code \\
\texttt{02-generate-parameter-sensitivity-original.R} & Figure 5: original code \\
\bottomrule
\end{tabular}
\end{table}

\subsection{Extracted Code by Commit}

Code from each historical commit extracted for comparison.
Following the 2026-04-08 reorganization, these snapshots live
in \texttt{analysis/archive/}:

\begin{verbatim}
analysis/archive/commit_42ac030/   # Initial (publication)
analysis/archive/commit_8609f12/   # Ron Thomas edits
analysis/archive/commit_f6ee86d/   # Autocorrelation typo fix
analysis/archive/commit_f70f86d/   # Carryover added
analysis/archive/commit_325314f/   # Modified Db (broken)
analysis/archive/original_R/       # Copy of 42ac030
\end{verbatim}

The earlier exploratory tidyverse pipeline (\texttt{analysis/2025/})
has been moved to an external archive at
\texttt{\textasciitilde/Dropbox/prj/alz/10-pmsimstats-ng/archive/2025/}
and is no longer part of the working repository. Its successor
is the maintained tidyverse implementation in
\texttt{implementations/tidyverse/} (see Section~\ref{sec:impl}).

\section{Documentation (\texttt{docs/})}

The full documentation index is in
\texttt{00-documentation-index.tex}; the table below summarizes
each document by number. Sources are in \texttt{docs/}; rendered
PDFs are in \texttt{docs/pdf/}.

\begin{table}[h]
\centering
\small
\begin{tabular}{rll}
\toprule
\# & Document & Content \\
\midrule
01 & Codebase Overview & Package structure, parameters, dependencies (this doc) \\
02 & Two DGP Architectures & White paper: Architecture~A vs.\ B power under carryover \\
03 & Audit and Revision Report & Seven DGP corrections, PD resolution, deliverables \\
04 & Revised Power Analysis & Architecture~B problems, fixes, validated results \\
06 & AR(1) Residual Correlation & \texttt{lmer} to \texttt{nlme} transition with Type I validation \\
07 & Positive Definiteness Failures & PD failure prevalence and causes (Arch~B) \\
08 & BM-BR Correlation Decay & Pharmacokinetic derivation of $\lambda_{\text{cor}}$ (Arch~B) \\
09 & Carryover Correlation Artifact & BM-BR correlation artifact analysis (Arch~B) \\
10 & Carryover Analysis Assessment & Carryover characterization in actual data \\
13 & Figure 4 Walkthrough & Pre-audit methodology reference (historical) \\
14 & Trial Analysis Guide & Collaborator goal implementation \\
15 & Tidyverse Alignment Plan & Planning record (largely executed in April 8 merge) \\
16 & Tidyverse Code Review & Review of pre-merge \texttt{analysis/2025/} code (historical) \\
17 & Simplification Plan & Exploratory; not adopted in April 8 refactor \\
18 & Response Parameter Sensitivity & Figure 5 sensitivity before/after (Arch~B) \\
19 & Mean Moderation Implementation & Architecture~A derivation (\texttt{.tex} canonical, \texttt{.md} summary) \\
20 & Annotated Bibliography & N-of-1 simulation methodology literature \\
21 & Architecture Comparison Summary & Brief empirical contrast of Arch~A vs.\ B \\
22 & Journal Target Recommendation & Publication strategy for white paper \texttt{02} \\
\bottomrule
\end{tabular}
\end{table}

\section{Bundled Data (\texttt{data/})}

\begin{table}[h]
\centering
\begin{tabular}{ll}
\toprule
File & Content \\
\midrule
\texttt{results\_core.rda} & Published Figure 4 results (initial commit, 500 reps) \\
\texttt{extracted\_bp.rda} & Baseline parameters (BP mean/SD, CAPS mean/SD) \\
\texttt{extracted\_rp.rda} & Response parameters (Gompertz max/disp/rate/sd) \\
\texttt{CTdata.rda} & Actual clinical trial data \\
\texttt{results\_maxes.rda} & Figure 5A results (response maxima) \\
\texttt{results\_rates.rda} & Figure 5B results (response rates) \\
\texttt{results\_trajectories.rda} & Factor trajectory data \\
\bottomrule
\end{tabular}
\end{table}

\section{Key Parameters}

\subsection{Current Configuration (Revised Code)}

\begin{table}[h]
\centering
\begin{tabular}{lll}
\toprule
Parameter & Value & Rationale \\
\midrule
\texttt{dgp\_architecture} & \texttt{"mvn"} (default) & Architecture~B; \texttt{"mean\_moderation"} for A \\
Autocorrelation & AR(1), $\rho = 0.7$ & PD-feasible, clinically realistic \\
\texttt{c.bm} (Arch~B) & $\{0, 0.25, 0.45\}$ & Max PD-feasible $\approx 0.45$ (OL+BDC) \\
\texttt{c.bm} (Arch~A) & up to $0.95$ & PD constraint absent under Arch~A \\
Carryover $t_{1/2}$ & $\{0, 0.5, 1.0\}$ weeks & Clinically realistic \\
\texttt{lambda\_cor} & $\ln 2 / t_{1/2}$ (auto) & Pharmacokinetic derivation \\
Cross-factor (same $t$) & \texttt{c.cf1t} $= 0.2$ & Unchanged from publication \\
Cross-factor (diff $t$) & \texttt{c.cfct} $\cdot \rho^{|t_i-t_j|}$ & AR(1) decay applied \\
Analysis model & \texttt{nlme::lme} + \texttt{corCAR1} & Required for AR(1) Type I control \\
PD corrections & 0/162 (0\%) & All matrices valid \\
\bottomrule
\end{tabular}
\end{table}

\subsection{Publication Configuration (Original Code)}

\begin{table}[h]
\centering
\begin{tabular}{lll}
\toprule
Parameter & Value & Issue \\
\midrule
Autocorrelation & CS, rho = 0.8 & PD failures at c.bm > 0.25 \\
c.bm & \{0, 0.3, 0.6\} & 0.6 exceeds PD-feasible range \\
Carryover t\_half & \{0, 0.1, 0.2\} weeks & Too short for genuine effects \\
BM-BR correlation & Step function & Power drop artifact \\
Analysis model & lmer (no corCAR1) & Correct under CS \\
PD corrections & 40/162 (24.7\%) & Silent, uncontrolled \\
\bottomrule
\end{tabular}
\end{table}

\section{Dependencies}

Core: \texttt{data.table}, \texttt{nlme}, \texttt{lme4}, \texttt{lmerTest}, \texttt{corpcor},
\texttt{MASS}, \texttt{ggplot2}

Performance: \texttt{furrr}, \texttt{future} (parallel processing)

Legacy (in DESCRIPTION but used minimally): \texttt{svMisc},
\texttt{tictoc}, \texttt{ggpubr}, \texttt{gridExtra}

\vfill
\noindent\rule{\textwidth}{0.4pt}
{\footnotesize
Rendered on 2026-04-09 at 18:27 PDT.\\
Source: \textasciitilde/prj/alz/10-pmsimstats-ng/pmsimstats-ng/docs/01-codebase-overview.tex}
\end{document}
